\section{Conditional matching elimination and an initial-mask obstruction}
\label{sec:conditional-type-matching}

Quadrilateral compatibility suggests a polynomial preprocessing rule: test a
type after setting its two local competitors to zero, and delete it if matching
alone then forces it to vanish.  The rule is sound, but the retained six-source
experiment finds no deletion.  A local matching identity explains why a complete
LP support test cannot improve the initial linear test.  Exact original-coordinate
witnesses establish conditional feasibility for all $276$ initial types of these
sources, including the suffix not completed by the earlier sufficient-closure
experiment.  This is a theory and evidence checkpoint, with no native search
integration or claimed recognition speedup.

\subsection{A sequential rule with original-equation provenance}

Fix a finite source triangulation with $t$ tetrahedra and let $M$ be its
original STANDARD normal matching matrix, with four triangle and three
quadrilateral columns per tetrahedron.  Write $C$ for a Q matching matrix
obtained by eliminating triangle differences.  For a retained Q mask $J$, let
\[
 K_J=\{q\ge0:Cq=0,\quad q_b=0\text{ for }b\notin J\}.
\]
Every $q\in K_J$ has a nonnegative triangle lift: first take a real matching
lift, then add triangle vertex links until every triangle coordinate is
nonnegative.  The source is finite, so these are ordinary finite vectors.
Canonical lifting subtracts the minimum triangle coordinate independently
at each global vertex.

For a type $a=(\tau,i)\in J$, define the conditional mask
\[
 J_a=J\setminus\{(\tau,j):j\ne i\}.
\]
A compatible vector with $q_a>0$ has both competing Q coordinates zero and
therefore belongs to $K_{J_a}$.  It follows that
\[
 \bigl(q_a=0\text{ for every }q\in K_{J_a}\bigr)
 \quad\Longrightarrow\quad
 \text{deleting $a$ preserves every compatible vector in }K_J.
\]
This implication can be applied sequentially, starting a new conditional
test from the current mask after each accepted deletion.  Only the tested
type $a$ is licensed for global deletion.  A second type forced to vanish
under the same conditional hypothesis need not vanish for compatible
vectors using a different type at $\tau$.

A sufficient zero proof can be checked directly against the original
STANDARD rows.  Form an exact rational linear combination of those rows
and require every triangle coefficient to cancel.  On the current
conditional mask, require the remaining Q coefficients to be nonnegative,
except at coordinates already proved zero earlier in this same local
proof.  Their nonnegative weighted sum vanishes, so every coordinate
with strictly positive coefficient also vanishes.  Coefficients outside
the conditional mask may have either sign because those coordinates are
zero.  The local forced set must start empty for every new conditional
query.  These checks establish the sequential inference without trusting
the producer's elimination or sign-closure computation.

With initially at most $3t$ types, a sorted scan restarting after every
deletion makes $O(t^2)$ queries.  The one- and two-row sufficient closure
used in the archived probe is polynomial, but its general fixed point is
not a complete nonnegative-feasibility certificate.  In particular, a
capped query provides no inference about whether its tested coordinate
can be positive.  The retained probe uses private arbitrary-mask algebra;
its raw provenance does not broaden the compatible-only meaning of an
existing native certificate schema.

\subsection{Mask monotonicity does not change the full-source parameter}

For nested masks $J'\subseteq J$, the cone $K_{J'}$ is a coordinate face
of $K_J$.  By Lemma~\ref{lem:transversal-face-rays}, its extreme rays are
extreme rays of $K_J$.  Zero-extension does not change the source, the
canonical formal Euler function $F$, or a ray's complete conflict set.
Consequently the hypergraph of incompatible positive extreme rays on
$J'$ is a subhypergraph of the corresponding hypergraph on $J$, and
\[
                              h_{J'}\le h_J.
\]
As in Section~\ref{sec:euler-transversal}, compatible positive rays are
excluded from this conflict hypergraph, and an empty hypergraph has
transversal number zero.

This is monotonicity of a \emph{mask parameter}.  The full-source $h(T)$
of the unchanged triangulation has not changed.  A complete negative
tree on the final mask, splitting only in $B$, proves $h_J\le |B|$ for
that mask.  Composing its negative conclusion with sound conditional
deletions can prove that the original admissible cone is nonpositive;
it does not establish $h(T)\le |B|$ for the full formal cone.  No claim
about transporting $h$ through a topological simplification is involved.

\subsection{A local block-ones matching identity}

For a tetrahedron $s$, define a signed STANDARD vector $z_s$.
Its four triangle coefficients in $s$ are all $-1$, and its three Q
coefficients there are all $1$; all other coordinates are zero.  On each face and
at each normal arc type, exactly one triangle and one quadrilateral
contribute.  Their sum is zero.  Thus $z_s$ has zero trace on every
face separately, and
\[
                                  Mz_s=0.
\]
The argument includes identified faces and self-attachments: both
sides of any matching equation see zero trace.  Its Q projection $r_s$
is the vector equal to one on the three Q coordinates of $s$ and zero
elsewhere, and hence $Cr_s=0$.

The signed vector $z_s$ is not a nonnegative normal surface.  This
matching identity alone says nothing about its Euler characteristic
or geometric realization.  Here it is used only to repair Q signs in
a linear matching solution.

\begin{proposition}[Complete initial conditional feasibility test]
\label{prop:conditional-initial-linear}
Let $J$ be the full set of $3t$ Q types and $a=(\tau,i)$.  The following
are equivalent:
\begin{enumerate}
\item There is a nonnegative STANDARD matching vector supported in Q
      mask $J_a$ whose $a$ coordinate is positive.
\item The $a$ coordinate is not identically zero on the linear Q
      matching space supported in $J_a$.
\end{enumerate}
When they hold, an integral witness can be constructed using exact
linear algebra and the vectors $r_s$, without an LP.
\end{proposition}
\begin{proof}
Projection proves the first implication.  Conversely choose a rational
linear matching vector with nonzero $a$ coordinate, dividing by that
coordinate so that $q_a=1$.  For each $s\ne\tau$, set
\[
 \lambda_s=\max\{0,-\min(q_{s,0},q_{s,1},q_{s,2})\}
 \quad\hbox{and replace }q\hbox{ by }q+\lambda_s r_s.
\]
Every such block retains all three Q types, so this operation stays
within $J_a$, preserves matching, and leaves $q_a=1$ and its competitors
zero.  After all replacements every Q coordinate is nonnegative.
Take a rational triangle lift.  Subtracting its minimum triangle
coordinate at each global vertex adds a triangle vertex-link matching
vector and makes every triangle coordinate nonnegative.  Finally
clear denominators, optionally divide the common gcd, and obtain the
claimed nonnegative integral STANDARD witness.
\end{proof}

The proof does not require a one-vertex source or a knot-type hypothesis.
It uses the ordinary finite matching system and its triangle vertex-link
kernel.  The retained sources happen to be one-vertex compact torus-boundary
sources authenticated from diagrams.

If the target coordinate vanishes on every nullspace basis vector,
its coordinate functional belongs to the matching row space.  In
reduced row echelon form, its pivot equation has no free-variable
terms and is precisely $q_a=0$.  The one-row forcing pass already
detects that unit equation.  Thus a complete LP support test at the
full initial mask cannot find a target zero missed by this linear
criterion and its unit-row check.  This is a completeness statement
for the \emph{initial conditional target test}, not for arbitrary
row-sign closure on arbitrary retained masks.

The construction has polynomial bit complexity.  There are $O(t)$
original matching rows and columns with bounded integer entries;
exact elimination coefficients are ratios of minors of these matrices.
Determinant bounds give polynomial bit lengths.  The block repairs,
triangle lift and vertex shifts use polynomially many rational
operations.  Clearing the $O(t)$ denominators and reducing their gcd
also gives a vector of polynomial bit length.  The explicit callback
counts below are implementation-work records, not this bit-operation
bound and not elapsed-time measurements.

\subsection{Why later masks need a different argument}

The block repair relies on all three types being retained in every
tetrahedron other than $\tau$.  After a deletion, a needed $r_s$ may
lie outside the retained mask.  The following exact algebra example
shows why the local block identities alone do not extend
Proposition~\ref{prop:conditional-initial-linear} to all later masks.

Use two blocks $(a,b,c)$ and $(d,e,f)$ with equations
\[
                a-b=0,\qquad -a-b+2c+d-e=0.
\]
Each equation has coefficient sum zero in each block.  Both local
block-ones vectors therefore belong to the kernel.  The initial
single-type tests at $a$ and $b$ force them to zero by linear unit
equations, so a sequential rule can delete both.  On the later
conditional mask for $d$, only $c,d$ remain and the equation is
$2c+d=0$.  The linear kernel has vectors with $d\ne0$, for example
$c=-d/2$, but nonnegative coordinates force $d=0$.
This is an abstract matching-algebra counterexample.  It is not
asserted to be realizable by a triangulated manifold or to be a
geometric fixture of the implementation.

\subsection{Two retained experiments on identical authenticated sources}

The archive \path{fast/conditional_type_research} preserves both
experiments byte-for-byte, including generation scripts, all returned
and partial data, both old replay results and both original manifests.
The labelled sources are exactly the enabled fixed-U1 endpoints from
Section~\ref{sec:pachner-descent-crushing}.  Their original PDs and
mixed source chains are copied in the exact staged provenance archive.
The labels $T,F,TT,T25,T27,W34$ denote the trefoil, figure-eight, visible
trefoil sum, and the braid inputs $(2,[1]^5)$, $(2,[1]^7)$ and
$(3,[1,-2]^4)$, respectively.  No primeness assertion about $W34$ is
needed for this source experiment.

\begin{center}
\begin{tabular}{lrrrrr}
Source & $t$ & Closure attempts/returns & Closure checks & Witnesses & Witness checks\\
\hline
$T$   & 7  & $21/21$ & $7\,924$       & 21 & $3\,047$\\
$F$   & 10 & $30/30$ & $30\,890$      & 30 & $6\,242$\\
$TT$  & 16 & $48/48$ & $212\,610$     & 48 & $16\,034$\\
$T25$ & 11 & $33/33$ & $52\,604$      & 33 & $7\,559$\\
$T27$ & 17 & $51/51$ & $290\,842$     & 51 & $18\,107$\\
$W34$ & 31 & $33/32$ & $1\,000\,000$ & 93 & $60\,359$\\
\hline
Total &    & $216/215$ & $1\,594\,870$ & 276 & $111\,348$
\end{tabular}
\end{center}

Each experiment allowed one million explicit callbacks per source,
six million total, with no wall-clock cutoff.  The first probe used
the sufficient one- and two-row closure.  Five sources completed
their scans, while $W34$ stopped at the conditional test $(10,2)$.
All $215$ returned queries supplied no local forcing packet; there
were zero global deletions.  There was therefore no genuine returned
packet with another locally forced type on which to run an
other-forced-global-deletion corruption control.  This absence is
retained explicitly, rather than filled with a synthetic geometry
claim.  Its archived replay validates source binding, scan transitions
and any claimed inference, but does not rerun closure or certify its
recorded kernel ranks.

The separate witness construction applies
Proposition~\ref{prop:conditional-initial-linear}.  All $276$ tests
complete, with no cap, error or zero-linear-target diagnostic.  Every
retained vector is checked independently for literal nonnegative
integer STANDARD coordinates, all original matching equations, a
positive tested Q coordinate and zero competing coordinates at that
tetrahedron.  The checker reconstructs each conditional mask and
replays the original diagram-bound source chains.  It also rejects
$12$ targeted corruptions, including an original matching residue,
boolean coordinate or type substitutions, and a positive competing
Q coordinate.

This second experiment establishes all-type initial conditional
feasibility on all six sources.  It covers the $W34$ suffix left
untested by the first probe, without retrospectively calling that
earlier capped scan complete.  The witnesses may have incompatible
Q coordinates in other tetrahedra.  They provide no assertion about
Euler sign, extremality, embedded normal surfaces or recognition
outcomes.  In particular, no complete LP run is needed to refute
initial matching-zero elimination on these sources, and none was
performed in either experiment.

The portable replay invokes no kernel constructor, sign closure, LP,
ray enumeration or topology discovery.  It checks exact archived
bytes against both stored and uncompressed hashes and records fresh
hashes of its current native, package and article dependencies.
Embedded generation pins remain historical.  Recorded construction
counters receive type, range and continuity checks; their operation
traces are not independently reproduced.  Independent replay is
excluded from those counters, and there is no timing or speedup claim.
These conditional-matching experiments add no native elimination rule.
These results close this
particular matching-only shortcut on the tested sources; they do not
bound general $h(T)$ or establish quasi-polynomial recognition.
